%============================================================
%  Bd. 16 - Kommutative Halbgruppen
%============================================================

\documentclass[main.tex]{subfiles}

\ifSubfilesClassLoaded{
  \BandSetupStandalone{B16}
}{
}

\title{Bd. 16 - Kommutative Halbgruppen}
\author{Martin Kunze}
\date{}
\setcounter{file}{16}

\begin{document}

\title{Bd. 16 - Kommutative Halbgruppen}
\author{Martin Kunze}
\date{}
\hypersetup{pageanchor=false}
\maketitle
\hypersetup{pageanchor=true}
\setcounter{tocdepth}{3}
\tableofcontents
\setcounter{file}{16}
\renewcommand{\FormulaBandID}{16}

\chapter{Einführung}

Dieser Band untersucht Halbgruppen mit kommutativer Operation.

\chapter{Kommutative Halbgruppen}

\section{Definition}

\FormulaDefDeltaK[Begriff der kommutativen Halbgruppe]{%
  \CommutativeSemigroupAlg{A}{\star}%
}{CommutativeSemigroupDef}{%
  \DeltaRow{Mengen}{A\dsep x\dsep y}
  \DeltaRow{\textbf{Axiome}}{}
  \DeltaPrem{Halbgruppe}{\SemigroupAlg{A}{\star}}
  \DeltaRow{Kommutativität}{%
    x,y\in A\vdash x\star y=y\star x}
  \DeltaRow{\textbf{Neues Symbol}}{}
  \DeltaPrem{Kommutative Halbgruppe}{%
    \CommutativeSemigroupAlg{A}{\star}}
}

\FormulaAxiomDeltaK[Zugriff auf die Halbgruppenstruktur]{%
  \CommutativeSemigroupAlg{A}{\star}
  \vdash\SemigroupAlg{A}{\star}%
}{CommutativeSemigroupBaseAxiom}{%
  \DeltaRow{Mengen}{A}
}

\FormulaAxiomDeltaK[Zugriff auf die Kommutativität]{%
  \CommutativeSemigroupAlg{A}{\star}\dsep x,y\in A
  \vdash x\star y=y\star x%
}{CommutativeSemigroupCommutativityAxiom}{%
  \DeltaRow{Mengen}{A\dsep x\dsep y}
}

\FormulaThmDeltaK[Dreigliedriges Vertauschungsgesetz]{%
  \CommutativeSemigroupAlg{A}{\star}\dsep a,b,c\in A
  \vdash (a\star b)\star c=(a\star c)\star b%
}{CommutativeSemigroupThreeTermExchange}{%
  \DeltaRow{Mengen}{A\dsep a\dsep b\dsep c}
}
\begin{tabproofwide}
  \proofstepwidestar[1]{\CommutativeSemigroupAlg{A}{\star}}{\rA}
  \proofstepwidestar[2]{a\in A}{\rA}
  \proofstepwidestar[3]{b\in A}{\rA}
  \proofstepwidestar[4]{c\in A}{\rA}
  \proofstepwidestar[1]{\SemigroupAlg{A}{\star}}{%
    \FormulaRefAuto{CommutativeSemigroupBaseAxiom}[ax]{1}}
  \proofstepwidestar[1,2,3,4]{%
    (a\star c)\star b=a\star(c\star b)}{%
    \FormulaRefAuto{SemigroupAssociativityAxiom}[ax]{5,2,4,3}}
  \proofstepwide[1,2,3,4]{(a\star b)\star c}{=}{a\star(b\star c)}{%
    \FormulaRefAuto{SemigroupAssociativityAxiom}[ax]{5,2,3,4}}
  \proofstepwide[1,3,4]{}{=}{a\star(c\star b)}{%
    \FormulaRefAuto{CommutativeSemigroupCommutativityAxiom}[ax]{1,3,4}}
  \proofstepwide[1,2,3,4]{}{=}{(a\star c)\star b}{%
    \FormulaRefAuto{a=b\vdash b=a}{6}}
  \proofstepwidestar[1,2,3,4]{%
    (a\star b)\star c=(a\star c)\star b}{%
    \rChain{7,9}}
\end{tabproofwide}

\FormulaThmDeltaK[Viergliedriges Vertauschungsgesetz]{%
  \begin{aligned}[t]
    &\CommutativeSemigroupAlg{A}{\star}\dsep a,b,c,d\in A\vdash{}\\[-2pt]
    &(a\star b)\star(c\star d)=(a\star c)\star(b\star d)
  \end{aligned}
}{CommutativeSemigroupFourTermExchange}{%
  \DeltaRow{Mengen}{A\dsep a\dsep b\dsep c\dsep d}
}
\begin{tabproofwide}
  \proofstepwidestar[1]{\CommutativeSemigroupAlg{A}{\star}}{\rA}
  \proofstepwidestar[2]{a\in A}{\rA}
  \proofstepwidestar[3]{b\in A}{\rA}
  \proofstepwidestar[4]{c\in A}{\rA}
  \proofstepwidestar[5]{d\in A}{\rA}
  \proofstepwidestar[1]{\SemigroupAlg{A}{\star}}{%
    \FormulaRefAuto{CommutativeSemigroupBaseAxiom}[ax]{1}}
  \proofstepwidestar[1,2,3]{a\star b\in A}{%
    \FormulaRefAuto{SemigroupClosureAxiom}[ax]{6,2,3}}
  \proofstepwidestar[1,2,4]{a\star c\in A}{%
    \FormulaRefAuto{SemigroupClosureAxiom}[ax]{6,2,4}}
  \proofstepwidestar[1,2,3,4,5]{%
    \begin{aligned}[t]
      &(a\star b)\star(c\star d)\mathrel{=}\\[-2pt]
      &\qquad\bigl((a\star b)\star c\bigr)\star d
    \end{aligned}}{%
    \FormulaRefAuto{a=b\vdash b=a}{%
      \FormulaRefAuto{SemigroupAssociativityAxiom}[ax]{6,7,4,5}}}
  \proofstepwidestar[1,2,3,4,5]{%
    \begin{aligned}[t]
      &\bigl((a\star b)\star c\bigr)\star d\mathrel{=}\\[-2pt]
      &\qquad\bigl((a\star c)\star b\bigr)\star d
    \end{aligned}}{%
    \FormulaRefAuto{CommutativeSemigroupThreeTermExchange}{1,2,3,4}}
  \proofstepwidestar[1,2,3,4,5]{%
    \begin{aligned}[t]
      &\bigl((a\star c)\star b\bigr)\star d\mathrel{=}\\[-2pt]
      &\qquad(a\star c)\star(b\star d)
    \end{aligned}}{%
    \FormulaRefAuto{SemigroupAssociativityAxiom}[ax]{6,8,3,5}}
  \proofstepwidestar[1,2,3,4,5]{%
    (a\star b)\star(c\star d)=(a\star c)\star(b\star d)}{%
    \rChain{9,11}}
\end{tabproofwide}

\section{Beispiele}

\FormulaThmDeltaK[Die Addition bildet eine kommutative Halbgruppe]{%
  \CommutativeSemigroupAlg{\mathbb{N}}{+}%
}{NaturalAdditionCommutativeSemigroup}{%
  \DeltaRow{Mengen}{\mathbb{N}\dsep a\dsep b}
}
\begin{tabproof}
  \proofstep{}{\SemigroupAlg{\mathbb{N}}{+}}{%
    \FormulaRefAuto{NaturalAdditionSemigroup}[thm]}
  \proofstep{2}{a\in\mathbb{N}}{\rA}
  \proofstep{3}{b\in\mathbb{N}}{\rA}
  \proofstep{2,3}{a+b=b+a}{%
    \FormulaRefAuto{PeanoAddCommutative}[thm]{2,3}}
  \proofstep{}{\CommutativeSemigroupAlg{\mathbb{N}}{+}}{%
    \FormulaRefAuto{CommutativeSemigroupDef}[def]{1,4}}
\end{tabproof}

\FormulaThmDeltaK[Die Vereinigung bildet auf einer Potenzmenge eine
kommutative Halbgruppe]{%
  \CommutativeSemigroupAlg{\powerset(M)}{\cup}%
}{PowerSetUnionCommutativeSemigroup}{%
  \DeltaRow{Mengen}{M\dsep X\dsep Y\dsep Z}
}
\begin{tabproof}
  \proofstep{1}{X\in\powerset(M)}{\rA}
  \proofstep{2}{Y\in\powerset(M)}{\rA}
  \proofstep{1}{X\subseteq M}{%
    \FormulaRefAuto{x \in \powerset(A)\;\eqvdash\; x \subseteq A}{1}}
  \proofstep{2}{Y\subseteq M}{%
    \FormulaRefAuto{x \in \powerset(A)\;\eqvdash\; x \subseteq A}{2}}
  \proofstep{1,2}{X\cup Y\subseteq M}{%
    \FormulaRefAuto{A \subseteq C,\, B \subseteq C \vdash A \cup B \subseteq C}{3,4}}
  \proofstep{1,2}{X\cup Y\in\powerset(M)}{%
    \FormulaRefAuto{x \in \powerset(A)\;\eqvdash\; x \subseteq A}{5}}
  \proofstep{7}{Z\in\powerset(M)}{\rA}
  \proofstep{}{(X\cup Y)\cup Z=X\cup(Y\cup Z)}{%
    \FormulaRefAuto{(A \cup B) \cup C = A \cup (B \cup C)}}
  \proofstep{}{\SemigroupAlg{\powerset(M)}{\cup}}{%
    \FormulaRefAuto{SemigroupDef}[def]{6,8}}
  \proofstep{}{X\cup Y=Y\cup X}{%
    \FormulaRefAuto{A \cup B = B \cup A}}
  \proofstep{}{\CommutativeSemigroupAlg{\powerset(M)}{\cup}}{%
    \FormulaRefAuto{CommutativeSemigroupDef}[def]{9,10}}
\end{tabproof}

\section{Primelemente}

Im nichtkommutativen Fall hängt ein Primbegriff von der gewählten linken,
rechten oder zweiseitigen Teilbarkeit ab.  Ohne weiteren Zusatz verwenden
wir \emph{Primelement} deshalb nur für kommutative Halbgruppen.  An die Stelle
von „keine Einheit“ tritt die auch ohne neutrales Element sinnvolle
Forderung, dass \(p\) nicht jedes Element teilt.  Damit legen wir als
Konvention fest, dass das von \(p\) erzeugte Hauptideal
\(A^1pA^1\) echt sein muss.

\FormulaDefDeltaK[Primelement einer kommutativen Halbgruppe]{%
  \SemigroupPrime{A,\star}{p}%
}{CommutativeSemigroupPrimeDef}{%
  \DeltaRow{Mengen}{A\dsep a\dsep b\dsep p\dsep q}
  \DeltaRow{Funktionen}{\star}
  \DeltaPrem{Kommutative Halbgruppe}{%
    \CommutativeSemigroupAlg{A}{\star}}
  \DeltaRow{Element}{p\in A}
  \DeltaRow{Echtheit}{%
    \exists q\in A\,
      \neg\SemigroupDivides{A,\star}{p}{q}}
  \DeltaRow{Primeigenschaft}{%
    \begin{aligned}[t]
      &a,b\in A\dsep
        \SemigroupDivides{A,\star}{p}{a\star b}\vdash{}\\[-2pt]
      &\SemigroupDivides{A,\star}{p}{a}\lor
       \SemigroupDivides{A,\star}{p}{b}
    \end{aligned}}
  \DeltaRow{\textbf{Neues Symbol}}{}
  \DeltaPrem{Primelement}{\SemigroupPrime{A,\star}{p}}
}

\paragraph{Beispiel.}
Die Halbgruppe
\[
  M=\{n\in\mathbb N\mid 2\leq n\}
\]
mit der Multiplikation besitzt kein neutrales Element.  Nach der vorstehenden
Definition sind ihre Primelemente dennoch genau die gewöhnlichen Primzahlen:
Für eine Primzahl \(p\) bezeugt \(p+1\in M\) die Echtheitsbedingung, denn
\(p\) teilt \(p+1\) nicht; die Primeigenschaft folgt aus dem Lemma von
Euklid.  Ist dagegen \(p=a\cdot b\) mit \(a,b\geq2\), so gilt
\(a,b<p\).  Daher teilt \(p\) weder \(a\) noch \(b\), wohl aber
\(a\cdot b=p\); somit verletzt jedes zusammengesetzte \(p\) die
Primeigenschaft.

\section{Morphismen}

\FormulaDefDeltaK[Homomorphismus kommutativer Halbgruppen]{%
  \CommutativeSemigroupHom{f}{A,\star}{B,\diamond}%
}{CommutativeSemigroupHomDef}{%
  \DeltaRow{Mengen}{A\dsep B}
  \DeltaRow{Funktionen}{f\dsep\star\dsep\diamond}
  \DeltaRow{\textbf{Axiome}}{}
  \DeltaPrem{Quellstruktur}{\CommutativeSemigroupAlg{A}{\star}}
  \DeltaPrem{Zielstruktur}{\CommutativeSemigroupAlg{B}{\diamond}}
  \DeltaPrem{Halbgruppenhomomorphismus}{%
    \SemigroupHom{f}{A,\star}{B,\diamond}}
  \DeltaRow{\textbf{Neues Symbol}}{}
  \DeltaPrem{Homomorphismus}{%
    \CommutativeSemigroupHom{f}{A,\star}{B,\diamond}}
}

\FormulaAxiomDeltaK[Quellstruktur]{%
  \CommutativeSemigroupHom{f}{A,\star}{B,\diamond}
  \vdash\CommutativeSemigroupAlg{A}{\star}%
}{CommutativeSemigroupHomSourceAxiom}{%
  \DeltaRow{Mengen}{A\dsep B}
  \DeltaRow{Funktionen}{f\dsep\star\dsep\diamond}
}

\FormulaAxiomDeltaK[Zielstruktur]{%
  \CommutativeSemigroupHom{f}{A,\star}{B,\diamond}
  \vdash\CommutativeSemigroupAlg{B}{\diamond}%
}{CommutativeSemigroupHomTargetAxiom}{%
  \DeltaRow{Mengen}{A\dsep B}
  \DeltaRow{Funktionen}{f\dsep\star\dsep\diamond}
}

\FormulaAxiomDeltaK[Zugriff auf den Halbgruppenhomomorphismus]{%
  \CommutativeSemigroupHom{f}{A,\star}{B,\diamond}
  \vdash\SemigroupHom{f}{A,\star}{B,\diamond}%
}{CommutativeSemigroupHomBaseHomAxiom}{%
  \DeltaRow{Mengen}{A\dsep B}
  \DeltaRow{Funktionen}{f\dsep\star\dsep\diamond}
}

\FormulaDefDeltaK[Isomorphismus kommutativer Halbgruppen]{%
  \CommutativeSemigroupIso{f}{A,\star}{B,\diamond}%
}{CommutativeSemigroupIsoDef}{%
  \DeltaRow{Mengen}{A\dsep B}
  \DeltaRow{Funktionen}{f\dsep\star\dsep\diamond}
  \DeltaRow{\textbf{Axiome}}{}
  \DeltaPrem{Homomorphismus}{%
    \CommutativeSemigroupHom{f}{A,\star}{B,\diamond}}
  \DeltaPrem{Bijektivität}{f\colon A\bij B}
  \DeltaRow{\textbf{Neues Symbol}}{}
  \DeltaPrem{Isomorphismus}{%
    \CommutativeSemigroupIso{f}{A,\star}{B,\diamond}}
}

\FormulaAxiomDeltaK[Homomorphismuszugriff]{%
  \CommutativeSemigroupIso{f}{A,\star}{B,\diamond}
  \vdash\CommutativeSemigroupHom{f}{A,\star}{B,\diamond}%
}{CommutativeSemigroupIsoHomAxiom}{%
  \DeltaRow{Mengen}{A\dsep B}
  \DeltaRow{Funktionen}{f\dsep\star\dsep\diamond}
}

\FormulaAxiomDeltaK[Bijektivitätszugriff]{%
  \CommutativeSemigroupIso{f}{A,\star}{B,\diamond}
  \vdash f\colon A\bij B%
}{CommutativeSemigroupIsoBijectiveAxiom}{%
  \DeltaRow{Mengen}{A\dsep B}
  \DeltaRow{Funktionen}{f\dsep\star\dsep\diamond}
}

\end{document}
